\chapimg{fig/ch04.jpg}
\chapterlead{\subhd{After this chapter you can} reproduce the blocks of the Results tab from the daily series of Engine~2, close the water balance of a year, compute the sums of the next 14 days, read the pinned-baseline differences with the right colour, and compute the reliability statistics of the Charts tab over twenty seasons.}
\chapter{The Dashboard: how every number is produced}\label{ch:dash}
\markboth{Dashboard}{Dashboard}

The Dashboard recomputes from the engine every time the state, the date or the weather changes, so it adds no model of its own: every number is a read-out or a sum of the daily series $P_k,I_k,E_k,T_k,D_k$ of Engine~2 (Chapter~\ref{ch:e2}) and of the annual outputs of Engine~1. This chapter does the sums.

\section{Block 1: the working date}
On day index $k$ the tab shows $\mathrm{Rain}=P_k$, $\mathrm{Irr}=I_k$, $\mathrm{Evap}=E_k$, $\mathrm{Transp}=T_k$, $\mathrm{Drain}=D_k$ (mm d$^{-1}$), the sensor depletion $f$ and the trigger $f_T$ (irrigation fires when $f\ge f_T$).

\exbox{One day of the dashboard}{The loam at the trigger $-300$ cm (depletion $f_T=0.468$), the first synthetic year (seed 1), working date day 150. Read the first block.}\sol  From the daily series: rain $P=1.0$ mm; rain reaching the soil after interception $0.3$ mm; irrigation $I=0.0$ mm; soil evaporation $E=0.54$ mm; transpiration $T=2.20$ mm; drainage $D=0.0$ mm; canopy cover $c=0.48$; the sensor layer holds 0.924 of its capacity so its relative depletion is $f=1-0.924=0.076$, below the trigger $f_T=0.468$. Irrigation fires on a day when the sensor depletion has reached $f_T$ (Eq.~6.1).


\section{Block 2: the next 14 days}
The sums over the following $N=14$ days are $\Sigma X=\sum_{j=1}^{14}X_{k+j}$ for $X\in\{P,I,E,T,D\}$. The bar says whether the days come from the weather record (the analogue year) or from the 16-day forecast when online weather is in use.

\exbox{The next fortnight}{Sum days 151 to 164.}\sol  $\Sigma P=13.6$ mm, $\Sigma I=15.0$ mm, $\Sigma E=7.4$ mm, $\Sigma T=39.8$ mm, $\Sigma D=0.0$ mm. Evapotranspiration of the fortnight is $\Sigma E+\Sigma T=47.2$ mm, a mean of $3.37$ mm d$^{-1}$. Without rain or irrigation the fortnight would draw 18.6 mm from storage (the model gives the actual change below).


\section{Block 3: the whole year and its closure}
For the calendar year the tab gives rain, irrigation, other evaporation (interception, litter, ponded water, snow), soil evaporation, transpiration, drainage and the storage change; the balance closes by the identity
\begin{equation}\boxed{\Delta S=P+I-E_{\rm other}-E-T-D.}\end{equation}

\begin{table}[!htbp]\centering\small\caption{Block 3 of the Results tab.}\label{tab:year}
\begin{tabular}{lr}\toprule
\hd{Quantity} & \hd{mm (year 1, loam, $-300$ cm)}\\\midrule
Rain $P$ & 606.1\\
Irrigation $I$ & 90.0\\
Other evaporation $E_{\rm other}$ & 61.7\\
Soil evaporation $E$ & 158.5\\
Transpiration $T$ & 229.6\\
Drainage $D$ & 200.2\\
Storage change $\Delta S$ & 46.1\\
Deepest roots (cm) & 82\\\bottomrule\end{tabular}\end{table}

\exbox{Closing the year}{Verify the closure of Table~\ref{tab:year}.}\sol  $\Delta S=P+I-E_{\rm other}-E-T-D=606.1+90.0-61.7-158.5-229.6-200.2=46.1$ mm. Independently, the root-zone store of the 100 cm profile rose from 122.4 mm on 1 January to 168.5 mm on 31 December, a change of 46.1 mm. The two agree because $E_{\rm other}$ is, by construction, the rain that did not reach the soil: $E_{\rm other}=\sum(P_k-P_k^{\rm soil})=61.7$ mm.


\section{Block 4 and the Charts tab: twenty seasons}
Block 4 shows the 20-year averages of the annual model, and each chart has a line stating the mean, the 10--90\% range, the lowest and highest season and the variation between years. The variation is the coefficient of variation $\mathrm{CV}=100\,s/\bar x$ with $s$ the sample standard deviation.

\begin{table}[!htbp]\centering\small\caption{Reliability over twenty synthetic seasons (seeds 1--20), loam, trigger $-300$ cm, mm per year.}\label{tab:rel}
\begin{tabular}{lrrrrrr}\toprule
\hd{Quantity} & \hd{mean} & \hd{10\%} & \hd{90\%} & \hd{lowest} & \hd{highest} & \hd{CV (\%)}\\\midrule
Rain & 602 & 522 & 735 & 438 & 765 & 14\\
Irrigation & 99 & 58 & 137 & 45 & 150 & 32\\
Drainage & 207 & 134 & 283 & 111 & 359 & 31\\
Soil evaporation & 158 & 157 & 160 & 148 & 160 & 2\\
Transpiration & 230 & 229 & 230 & 228 & 230 & 0\\\bottomrule\end{tabular}\end{table}

\exbox{The statistics line}{From Table~\ref{tab:rel} write the line under the irrigation chart.}\sol  ``Irrigation: mean 99 mm; 10 to 90\% range 58 to 137 mm; lowest 45 mm, highest 150 mm; variation between years 32\%.'' The coefficient of variation is $100\times s/\bar x$ with $s=31.3$ mm and $\bar x=99$ mm. No season was rainfed in this sample.


\section{The pinned baseline}
Press Pin and the tiles of Block 4 are stored; afterwards each tile shows $\Delta_i=x_i-x_i^{\rm pin}$ with $\blacktriangle$ for an increase, $\blacktriangledown$ for a decrease, green when the change is in the better direction, red when worse and ``= no change'' when $|\Delta_i|<10^{-9}$. Lower is better for soil evaporation, drainage, irrigation and nitrogen leaching; higher is better for yield and irrigation efficiency; transpiration, evapotranspiration and rainfall are neutral.

\begin{table}[!htbp]\centering\small\caption{Pin at the trigger $-300$ cm, then move to $-1500$ cm (twenty-season means, mm per year).}\label{tab:pin}
\begin{tabular}{lrrrll}\toprule
\hd{Tile} & \hd{pinned (-300 cm)} & \hd{now (-1500 cm)} & \hd{$\Delta$} & \hd{arrow} & \hd{colour}\\\midrule
Soil evaporation & 157.8 & 156.1 & -1.7 & $\blacktriangledown$ & better (green)\\
Transpiration & 229.6 & 229.6 & 0.0 & = no change & neutral\\
Drainage & 206.9 & 156.2 & -50.7 & $\blacktriangledown$ & better (green)\\
Irrigation & 99.0 & 45.8 & -53.2 & $\blacktriangledown$ & better (green)\\
Rain & 602.0 & 602.0 & 0.0 & = no change & neutral\\\bottomrule\end{tabular}\end{table}

\exbox{Reading the pinned differences}{Read the tiles of Table~\ref{tab:pin}.}\sol  Irrigation falls from 99.0 to 45.8 mm (-53.2 mm): $\blacktriangledown$ and green, since less is better. Transpiration changes by 0.0 mm and is neutral (it carries no colour, because less transpiration may mean either less stress or a smaller crop). Drainage changes by -50.7 mm. A drier trigger therefore saves irrigation but the dashboard alone cannot say whether the saving is worth it: the yield tile (not computed here) must be read with it, as in Chapter~\ref{ch:ctrl}.


\section{The pump hydrograph of a design}
For drip and sprinkler designs the pump hydrograph is the step series $Q_{\rm pump}(t)=Q_z$ for $t\in[0,n_zt_{\rm run})$ and 0 afterwards: the zones or sets run in turn.

\exbox{A pump hydrograph}{For the sprinkler farm of Chapter~\ref{ch:sprinkler} ($n=3$ laterals at once, set time $T=371$ min) the farm has $N_m=18$ lateral positions, so $n_z=\lceil18/3\rceil=6$ sets. Find the step series.}\sol  The flow of one set is $Q_z=3\times12\times21.2$ L min$^{-1}=764$ L min$^{-1}=12.74$ L s$^{-1}$ (3 laterals of 12 sprinklers). $t_{\rm run}=371/60=6.18$ h. The pump runs at $Q_z$ for $n_zt_{\rm run}=37.1$ h, in steps of $6.18$ h, and the volume is $Q_z\times n_zt_{\rm run}\times3.6=1{,}701$ m$^3$ (check: gross depth $\times$ area of the sets).


\section{Summary}
Every Dashboard number is a read-out, a sum, a mean or a difference of the daily series: Block 1 reads day $k$; Block 2 sums 14 days; Block 3 closes $\Delta S=P+I-E_{\rm other}-E-T-D$; Block 4 and the charts use twenty seasons; the pin shows $\Delta=x-x^{\rm pin}$ with a colour by direction; the pump hydrograph is a step series. Because all of it comes from the same objects, the dashboard, the report and the CSV exports agree digit for digit.
